\chapter{Debugging MERN Applications}

A MERN bug can live in the frontend, the network layer, the backend route,
the controller logic, or the database itself. Guessing at random wastes
time; checking a fixed sequence of points, in order, tells you exactly
which layer the problem is in before you write a single fix.

\section{The Core Principle}

Each layer in the request/response path can only pass along what the layer
before it produced. If the request body arriving at the controller is
empty, no amount of controller logic will fix it --- the bug is upstream,
likely in how the frontend built the request. Working through the
checklist in order means you stop at the first point where reality departs
from what you expect, instead of debugging correct code because the actual
bug is somewhere else entirely.

\section{The Debugging Checklist}

\begin{enumerate}
\item \textbf{Browser Console} --- any thrown JS errors or warnings?
\item \textbf{Network tab} --- was the request even sent?
\item \textbf{Request URL} --- is it the exact endpoint you think it is?
\item \textbf{HTTP method} --- \texttt{GET}/\texttt{POST}/\texttt{PUT}/\texttt{DELETE} as expected?
\item \textbf{Request body} --- does it contain the data you expect, in the right shape?
\item \textbf{Headers} --- is \texttt{Content-Type} (and the auth cookie) actually present?
\item \textbf{Status code} --- 2xx, or does it reveal a 4xx/5xx?
\item \textbf{Response body} --- what did the server actually say?
\item \textbf{Backend terminal} --- any stack trace or logged error?
\item \textbf{Route} --- is this request even reaching the route you expect?
\item \textbf{Controller} --- is the logic doing what you think with the input it received?
\item \textbf{Database} --- does the data actually look right in Compass/Atlas?
\end{enumerate}

\begin{diagrambox}[Debugging Checklist Flow]
\begin{verbatim}
 Browser Console
       |
       v
 Network tab (was request sent?)
       |
       v
 Request URL --> HTTP method --> Request body --> Headers
       |
       v
 Status code --> Response body
       |
       v
 Backend terminal (server-side errors/logs)
       |
       v
 Route --> Controller --> Database
       |
       v
 First mismatch found = the layer with the bug
\end{verbatim}
\end{diagrambox}

\begin{notebox}[NOTE]
The first six items (Console through Headers) are all frontend/network
layer. Status code and response body tell you whether the backend even ran
correctly. Everything from backend terminal onward is backend/database
layer. Most of the checklist's value is in quickly telling you which
\emph{half} of the stack to keep investigating.
\end{notebox}

\section{Worked Example: "Task Creation Silently Fails"}

The symptom: clicking "Create Task" does nothing visible --- no new task
appears, no error toast, nothing.

\begin{enumerate}
\item \textbf{Browser Console} --- no errors logged. Rules out a JS
exception in the click handler.
\item \textbf{Network tab} --- the \texttt{POST /api/tasks} request is
there. It was sent. Rules out the button not firing at all.
\item \textbf{Request URL} --- correct: \texttt{/api/tasks}.
\item \textbf{HTTP method} --- correct: \texttt{POST}.
\item \textbf{Request body} --- opening the request payload in DevTools
shows \texttt{\{\}} --- an empty object, even though the form clearly had a
title typed into it.
\item \textbf{Headers} --- checking the request headers shows
\texttt{Content-Type} is missing entirely; the request was sent with the
default browser content type instead of \texttt{application/json}.
\end{enumerate}

Found it: the axios call omitted the JSON content-type header, so Express's
\texttt{express.json()} body parser never parsed the payload, and
\texttt{req.body} landed in the controller as \texttt{\{\}}. The fix is on
the frontend, not the backend:

\begin{lstlisting}[language=JavaScript]
// Before -- api client missing default headers
const api = axios.create({ baseURL: import.meta.env.VITE_API_URL });

// After -- Content-Type set for every request
const api = axios.create({
  baseURL: import.meta.env.VITE_API_URL,
  headers: { "Content-Type": "application/json" },
  withCredentials: true,
});
\end{lstlisting}

\begin{warnbox}[WARNING]
Notice the bug was found at step 6 (Headers), well before reaching the
backend terminal, route, controller, or database. Jumping straight to
"add console.logs in the controller" would have wasted time debugging
correct backend code, since \texttt{req.body} being empty was a symptom of
a frontend request problem, not a controller bug.
\end{warnbox}

\begin{tipbox}[TIP]
Keep the browser's Network tab open with "Preserve log" enabled while
reproducing a bug. It's the single fastest way to confirm whether a
problem is frontend-side (request never sent, wrong URL/method/body) before
you even look at the backend.
\end{tipbox}

\whatwhyhow
{A fixed, ordered checklist -- console, network tab, URL, method, body, headers, status code, response body, backend terminal, route, controller, database -- isolates which layer of the stack a bug lives in.}
{Without a systematic order, it's easy to debug backend logic that was never actually reached because the real problem was a malformed or unsent frontend request.}
{Work the checklist top to bottom and stop at the first point where behavior departs from expectation -- that point is the layer containing the bug.}
